\chapter{Temporal Feature Integration and Hourly Loss Weighting}
\label{ch:time}

\section{Hypothesis Formulation and Temporal Encoding Objectives}

Prior experimental configurations provided models strictly with ordered cell identifier sequences, excluding explicit temporal timestamps. Because human mobility exhibits strong temporal periodicity (e.g., morning and evening commutes), excluding temporal context omits a key predictive feature.

The central research objective evaluated in this chapter is to investigate whether \emph{explicit temporal encoding (hour of day)} improves transition prediction accuracy, and whether loss weighting across specific transition hours forces the model to anticipate non-stationary mobility events.

\paragraph{Data Provenance.} Experiments evaluate active mobility dataset \Dmob{} (\textbf{Table~\ref{tab:datasets}}), comprising 400 training subscribers and \textbf{100 held-out evaluation subscribers}. Corresponding persistence baseline accuracy is 45.5\% on 80/20 evaluation splits and 43.1\% across whole-trajectory evaluation.

\subsection{Identification of Temporal Performance Degradation Patterns}

Individual subscriber error analysis revealed specific trajectories where model accuracy degraded sharply as context length increased.

% T5 -- Users whose accuracy collapses when given MORE context (image23, slide 38).
\begin{table}[t]
\centering\small
\caption[Users whose accuracy collapses when given more context]{The eight
users behind the aggregate of Table~\ref{tab:context-sweep}, scored with a
short context (30--60\,\% of the day) and then with a long one (80--90\,\%).
More history is supposed to help. For the bottom four it destroys the
prediction, which is what pointed at the \emph{end} of the day, and
therefore at time of day, as the missing signal. \smob, same trained model
evaluated at several context fractions.}
\label{tab:context-collapse}
\begin{tabular}{@{}lrrrl@{}}
\toprule
\textbf{User} & \textbf{Short ctx.} & \textbf{Long ctx.} & \textbf{$\Delta$} & \textbf{Reading} \\
\midrule
2046755 & 12.1\,\% & 10.8\,\% & $-1.3$  & \multirow{3}{*}{\begin{tabular}{@{}l@{}}\textbf{A}, difficulty is a per-user trait:\\ context changes nothing\end{tabular}} \\
599692  & 16.5\,\% & 16.7\,\% & $+0.2$  & \\
1243093 &  8.3\,\% & 15.6\,\% & $+7.3$  & \\
\addlinespace
1367716 & 60.4\,\% & 37.5\,\% & $-22.9$ & \textbf{B}, a real decline, still usable \\
\addlinespace
401222  & 32.2\,\% &  0.0\,\% & $-32.2$ & \multirow{4}{*}{\begin{tabular}{@{}l@{}}\textbf{C}, total collapse at the end\\ of the day: the model repeats the\\ last cell with up to 78\,\% confidence\end{tabular}} \\
1300439 & 40.7\,\% &  0.0\,\% & $-40.7$ & \\
1223357 & 33.7\,\% &  4.2\,\% & $-29.6$ & \\
1153529 & 20.0\,\% &  3.0\,\% & $-17.1$ & \\
\bottomrule
\end{tabular}

\vspace{3pt}
\begin{minipage}{\textwidth}\footnotesize
Accuracy at 30\,\% context correlates with accuracy at 80\,\% context at
$r = 0.77$ across all users, which is what makes group~A a \emph{trait} rather
than a context effect.
\end{minipage}
\end{table}


\textbf{Table~\ref{tab:context-collapse}} details eight representative subscriber trajectories exhibiting performance degradation at extended context fractions. The lower four entries demonstrate accuracy dropping to zero when evaluated on complete evening prefixes compared to initial morning segments. This empirical degradation indicates that trajectory context alone without temporal encoding fails to distinguish late-day transition probability shifts from mid-day patterns.

\section{Temporal Feature Representation and Experimental Hypotheses}

Temporal feature integration prepends dedicated hourly tokens (\cid{<H00>} through \cid{<H23>}) to cell event tokens within the sequence. Two distinct hypotheses were evaluated:

\begin{description}
  \item[Hypothesis 1 (Temporal Feature Sensitivity).] Providing explicit hour tokens enables the model to condition transition probabilities on diurnal temporal states (distinguishing morning vs. evening transition likelihoods).
  \item[Hypothesis 2 (Targeted Loss Weighting).] Assigning higher cross-entropy loss weights to specific \emph{transition windows}\footnote{\textbf{Transition window}: predefined temporal intervals where loss values are multiplied (e.g., $4\times$ multiplier for 18h--20h) to penalize transition prediction errors.} focuses gradient updates on high-mobility periods.
\end{description}

Initial loss weighting experiments assigned heuristic transition windows based on preliminary domain assumptions: \textbf{04:00--06:00} ($2\times$ weight) and \textbf{18:00--20:00} ($4\times$ weight).

\subsection{Ablation Analysis of Temporal Tokens}

To verify whether the model actively utilizes temporal tokens rather than ignoring them as sequence noise, an ablation experiment was performed.

% F9 -- The hour token: how it enters, and the ablation that tests whether it is read.
\begin{figure}[t]
\centering
% ---- (a) the 24-hour weight map ----------------------------------------
\fitwidth{%
\begin{tikzpicture}[font=\scriptsize,
  cell/.style={draw=clrule,fill=clsoft,minimum width=0.55cm,minimum height=0.50cm,inner sep=0pt},
  w2/.style={cell,draw=clgain,fill=clgain!25},
  w4/.style={cell,draw=clgain,fill=clgain!55},
]
\foreach \h in {0,1,2,3,6,7,8,9,10,11,12,13,14,15,16,17,20,21,22,23}
  {\node[cell] at (\h*0.58,0) {};}
\foreach \h in {4,5}   {\node[w2] at (\h*0.58,0) {$\times2$};}
\foreach \h in {18,19} {\node[w4] at (\h*0.58,0) {$\times4$};}
\foreach \h/\lbl in {0/00h,4/04h,6/06h,12/12h,18/18h,20/20h,23/23h}
  {\node[anchor=north,text=clrule] at (\h*0.58,-0.32) {\lbl};}
\node[anchor=west,font=\scriptsize,text=clrule,text width=13cm] at (0,0.72)
  {\textbf{(a)} loss weight applied to the target at each hour of the day. Bounds are
   half-open, so ``04--06h'' means hours 4 and 5.};
\end{tikzpicture}}

\vspace{0.45cm}
\replegend{\swatch{clmodel}~real hours \qquad \swatch{clnonsig}~hours shifted $+12$h}
\fitwidth{%
\begin{tikzpicture}
\begin{groupplot}[
  group style={group size=2 by 1, horizontal sep=1.6cm, ylabels at=edge left},
  repbar, width=0.43\linewidth, height=4.9cm,
  ylabel={top-1 accuracy (\%)}, xtick=data,
  every node near coord/.append style={font=\scriptsize},
]
% ---- (b) ablation ----
\nextgroupplot[title={\footnotesize (b) corrupting only the hour token},
  symbolic x coords={ha,hb,hc,hd},
  xticklabels={real,{shuffled\\\textcolor{clloss}{$-3.0\pt$}},{fixed 12h\\\textcolor{clloss}{$-0.7\pt$}},{shift $+12$h\\\textcolor{clgain}{$+0.5\pt$}}},
  x tick label style={font=\scriptsize,align=center},
  ymin=37, ymax=49.5, ytick={38,40,...,48}, /pgf/bar width=11pt]
\addplot[fill=clmodel,draw=none] coordinates {(ha,44.7)(hb,41.7)(hc,44.0)(hd,45.2)};
%
% ---- (c) where the shift lands ----
\nextgroupplot[title={\footnotesize (c) where the $+12$h shift lands},
  symbolic x coords={win,out},
  xticklabels={inside window,outside window},
  x tick label style={font=\scriptsize},
  ymin=36, ymax=51.5, ytick={36,38,...,50}, /pgf/bar width=13pt]
\addplot[fill=clmodel,draw=none] coordinates {(win,42.7)(out,45.3)};
\addplot[fill=clnonsig,draw=none] coordinates {(win,39.7)(out,47.1)};
\end{groupplot}
\end{tikzpicture}}
\caption{The hour-of-day ablation}
\label{fig:hour-ablation}
\figtext{%
Does the model actually read the hour? (a)~Each event enters the sequence as a
dedicated hour token followed by its \cid{cell\_id}, and the number in a
shaded box is the multiplier applied to the target's loss at that hour,
$\times2$ at 04--05h and $\times4$ at 18--19h. Bounds are half-open, so ``04--06h'' means hours 4 and 5, and multiplying the loss duplicates no data: it
only tells the optimiser that a mistake there costs more. (b)~The same trained
checkpoint on the same test users, with the cell sequence untouched and only
the hour token corrupted. The argument is the difference between two of the
bars: replacing every hour by a constant costs $0.7\pt$, while shuffling the
hours, so that the values are still real hours but in the wrong order, costs
$3.0\pt$. If the token merely occupied a slot, destroying its \emph{ordering}
could not cost more than removing it altogether; it costs four times more, so
the model has learned which hour goes with which cell. (c)~Shifting every hour
by twelve moves the aggregate by only $+0.5\pt$, and that near-zero hides a
swap: accuracy inside the transition windows falls $42.7 \rightarrow 39.7\,\%$
while accuracy outside them rises $45.3 \rightarrow 47.1\,\%$. The two halves
of the day exchanged their difficulty along with their labels. \smob{}, 80\,\%
context, $n = 1{,}747$ predictions.}
\end{figure}


\textbf{Figure~\ref{fig:hour-ablation}} illustrates the temporal weighting schedule in \textbf{\panelref{fig:hour-ablation}{a}}, the corruption strategies in \textbf{\panelref{fig:hour-ablation}{b}}, and compares model top-1 accuracy across three temporal token configurations in \textbf{\panelref{fig:hour-ablation}{c}} and \textbf{\panelref{fig:hour-ablation}{d}}: correct temporal tokens, constant dummy tokens (deleting temporal information), and randomly permuted temporal tokens (corrupting temporal ordering).

As shown in \textbf{\panelref{fig:hour-ablation}{c}}, replacing temporal tokens with a constant value reduces accuracy by 0.7 percentage points. Shuffling temporal tokens across event positions decreases accuracy by 3.0 percentage points ($4\times$ the drop of constant tokens) in \textbf{\panelref{fig:hour-ablation}{c}}. Furthermore, as detailed in \textbf{\panelref{fig:hour-ablation}{d}}, shifting hours by +12h swaps day and night performance, falling from 42.7 to 39.7 inside upweighted hours and rising from 45.3 to 47.1 outside them. This significant performance sensitivity under permuted temporal tokens confirms that the model actively conditions transition probability distributions on temporal token ordering, validating Hypothesis 1.

\subsection{Evaluation of Weighted Loss and Targeted Sampling}

To test Hypothesis 2, three joint modifications were implemented:
1. \textbf{Unequal Loss Weighting}: 18:00--20:00 transitions received a $4\times$ loss multiplier; 04:00--06:00 transitions received $2\times$.
2. \textbf{Targeted Sequence Sampling}: Sequences containing transition window events were sampled $3\times$ more frequently during batch generation, increasing transition batch representation from 45.1\% to 71.2\%.
3. \textbf{Full-Trajectory Evaluation Protocol}: Evaluation shifted from scoring exclusively the final 20\% prefix to scoring all trajectory event steps sequentially (reducing persistence baseline score from 45.5\% to 43.1\%).

\textbf{Figure~\ref{fig:targeted-sampling}} and \textbf{Table~\ref{tab:after-sampling}} evaluate the performance impact of these targeted optimization adjustments.

% F10 -- Targeted transition sampling: the windows are won, the rest is lost.
\begin{figure}[!htbp]
\centering
\replegend{\swatch{clbase}~the one-line rule \qquad \swatch{clmodel}~my model}
\fitwidth{%
\begin{tikzpicture}
\begin{axis}[repbar, width=\linewidth, height=5.6cm,
  symbolic x coords={wa,wb,wc,wd}, xtick=data,
  xticklabels={{04--06h\\128 moments},{18--20h\\604 moments},{both windows\\732 moments},{everywhere else\\7{,}893 moments}},
  /pgf/bar width=13pt,
  ymin=28, ymax=86, ytick={30,40,...,80},
  ylabel={correct answers per 100},
  x tick label style={font=\small,align=center},
]
\addplot[fill=clbase,draw=none]  coordinates {(wa,70.3)(wb,37.9)(wc,43.6)(wd,43.4)};
\addplot[fill=clmodel,draw=none] coordinates {(wa,71.1)(wb,40.6)(wc,45.9)(wd,40.8)};
\node[font=\scriptsize,text=clgain,anchor=south] at (axis cs:wa,78.5) {$+0.8\pt$};
\node[font=\scriptsize,text=clgain,anchor=south] at (axis cs:wb,48.5) {$+2.6\pt$};
\node[font=\scriptsize,text=clgain,anchor=south] at (axis cs:wc,53.5) {$+2.3\pt$};
\node[font=\scriptsize,text=clloss,anchor=south] at (axis cs:wd,51.5) {$-2.6\pt$};
\end{axis}
\end{tikzpicture}}
\caption{Emphasising two stretches of the day: the win, and the price}
\label{fig:targeted-sampling}
\figtext{%
\figpara{How to read the four pairs.}{Each pair is one group of graded moments,
with the number of moments printed under its label so the weight it carries is
visible. Grey-blue is the one-line rule, dark blue my model, both on exactly the
same moments, and the number above each pair is the distance between them. Before
this run, training slices containing a move inside one of the two chosen
stretches of the day were shown three times more often than the rest, and a
mistake at 18--20h was made to cost four times an ordinary one. Nothing else
moved.}

\figpara{The win is real; the price is larger.}{At 18--20h the model gains 2.6
correct answers per hundred, and 2.3 over the two chosen stretches combined. This
is the first time in the project that a trained model beat ``repeat the last
antenna'' on a segment chosen \emph{in advance} rather than found afterwards.
But the fourth pair is everything outside those stretches: 7{,}893 moments, more
than ten times the volume of the win, and the model loses 2.6 there. Making some
mistakes count more buys correct answers where the weight was put and pays for
them everywhere else, which is exactly what the mechanism is designed to do. The
trade is only worth making if the emphasised hours really are the hard ones, and
Table~\ref{tab:group-windows} shows they were not. \smob, 100 held-out people,
every position of the day graded; the 04--06h bar rests on 128 moments only and
should be read as ``not worse'' rather than as a win.}}

\end{figure}


% T6 -- What targeted transition sampling cost outside the windows (slide 46).
\begin{table}[!htbp]
\centering\footnotesize
\caption{The same change, measured everywhere else: seven cuts, seven losses}
\label{tab:after-sampling}
\begin{tabular}{@{}lrrrcc@{}}
\toprule
\textbf{Past given} & \textbf{Distance before} & \textbf{after} & \textbf{Change}
& \textbf{top-3, before $\rightarrow$ after} & \textbf{top-5, before $\rightarrow$ after} \\
\midrule
30\,\% & $-0.9$ & $-1.6$ & $-0.7$ & 61.4\,\% $\rightarrow$ 58.8\,\% & 69.1\,\% $\rightarrow$ 65.7\,\% \\
40\,\% & $-0.6$ & $-1.5$ & $-0.9$ & 61.7\,\% $\rightarrow$ 58.8\,\% & 69.0\,\% $\rightarrow$ 65.6\,\% \\
50\,\% & $-0.7$ & $-1.1$ & $-0.4$ & 61.4\,\% $\rightarrow$ 58.8\,\% & 68.6\,\% $\rightarrow$ 65.4\,\% \\
60\,\% & $-0.9$ & $-1.5$ & $-0.6$ & 61.7\,\% $\rightarrow$ 58.7\,\% & 68.7\,\% $\rightarrow$ 65.3\,\% \\
70\,\% & $-0.3$ & $-1.3$ & $-1.0$ & 61.4\,\% $\rightarrow$ 58.5\,\% & 68.6\,\% $\rightarrow$ 65.2\,\% \\
80\,\% \textit{(ref.)} & $-0.9$ & $-1.3$ & $-0.4$ & 62.3\,\% $\rightarrow$ 59.1\,\% & 70.2\,\% $\rightarrow$ 65.4\,\% \\
90\,\% & $-2.1$ & $-2.8$ & $-0.7$ & 64.6\,\% $\rightarrow$ 60.6\,\% & 72.3\,\% $\rightarrow$ 67.7\,\% \\
\bottomrule
\end{tabular}

\figsource{\nbllm}{train\_cellid\_llm.ipynb}{the targeted-sampling run, evaluated at seven cuts}
\end{table}


\subsection{Analysis of Empirical Window Misalignment}

Comparing the four bar pairs in \textbf{Figure~\ref{fig:targeted-sampling}} demonstrates that applying loss weighting and targeted sampling yields a +2.6 percentage point accuracy gain within the 18:00--20:00 window (the second bar pair of \textbf{Figure~\ref{fig:targeted-sampling}}) and +2.3 percentage points across both windows combined (the third bar pair of \textbf{Figure~\ref{fig:targeted-sampling}}), accompanied by a $-2.6$ percentage point accuracy decline across non-weighted temporal intervals (the fourth bar pair of \textbf{Figure~\ref{fig:targeted-sampling}}). \textbf{Table~\ref{tab:after-sampling}} confirms consistent behavior across context fractions.

Loss weighting effectively reallocates model optimization focus toward targeted temporal intervals. However, empirical evaluation revealed that the heuristic transition windows were not intrinsically harder than average: on full test trajectories, the persistence baseline scores 43.6\% \emph{inside} the guessed windows versus 43.4\% \emph{outside} them in \textbf{Figure~\ref{fig:targeted-sampling}}. Furthermore, early morning hours (04:00--06:00, the first bar pair of \textbf{Figure~\ref{fig:targeted-sampling}}) represent high regularity periods across all user cohorts (baseline accuracy 49\%--64\%). As documented by Song et al.\ \textbf{\cite{song2010limits}}, nocturnal hours exhibit maximum trajectory regularity, whereas transition difficulty dips around midday and early evening. Upweighting globally fixed hours applied uniformly across all subscribers failed to account for individual temporal mobility variations.

Because optimal transition windows vary across individual subscriber behaviors, global fixed temporal loss weighting is suboptimal. This insight motivates the user behavioral stratification approach presented in Chapter~\ref{ch:groups}.

\section{Conclusion}
\label{sec:retro-time}

This chapter validated the integration of explicit temporal tokens (\cid{<H00>}--\cid{<H23>}). Ablation analysis demonstrated that shuffling temporal tokens causes a 3.0 percentage point accuracy drop (\textbf{\panelref{fig:hour-ablation}{c}} and \textbf{\panelref{fig:hour-ablation}{d}}), confirming that model predictions depend on temporal sequence ordering.

Conversely, applying globally fixed hourly loss weighting (04:00--06:00 and 18:00--20:00) demonstrated that while loss weighting successfully forces local optimization (+2.6 percentage points at 18:00--20:00 in \textbf{Figure~\ref{fig:targeted-sampling}}), applying fixed hours across all subscribers penalizes non-targeted intervals. Because temporal mobility patterns differ significantly across subscribers, effective transition loss weighting requires individual user cohort stratification, as detailed in Chapter~\ref{ch:groups}.

